\chapter{Methods}

\section{Research Design}

The study is a fully crossed factorial experiment over retrieval configuration.
It crosses \thesisNumChunkSizes{} chunk sizes (\thesisChunkSizeSmall{} and
\thesisChunkSizeLarge{} tokens), \thesisNumStrategies{} retrieval strategies
(dense, BM25, and hybrid) and \thesisNumRepresentations{} indexing
representations (text-only and metadata-enriched), giving \thesisNumConditions{}
conditions. All \thesisNumQuestions{} benchmark questions ran through every
condition, producing \thesisNumGenerations{} generations. Implementation detail
for several sections is given in the appendix, which follows the order of this
chapter.

Differences between conditions are attributed to the three factors because the
following were held constant: the corpus, the questions, the prompt template,
the generator model, \thesisTopK{} retrieved chunks, the embedding model, the
random seed, and the decoding temperature. None varies in the grid.

\subsection{Hypotheses}

H1 to H3 are stated in Theory, each closing its own derivation. H1 is
non-directional. H2 is non-directional at the level tested, since the
design compares three specific strategies on this corpus, not the paradigms they
instantiate.
H3 is directional: company identity and period distinguish the right passage
from a near-identical wrong one, and are absent from the string the retriever
sees. Every test in this work is two-sided, including the test of H3, the one
directional hypothesis. A two-sided test spends part of its power on an outcome
the theory does not predict, so H3 is held to a stricter standard than it
requires.

Whether better retrieval yields more grounded, on-target answers is examined in
Results as an exploratory question rather than as a hypothesis. At these
discordant-pair counts it could not be tested: the exact test's smallest
attainable \emph{p} value sits above any conventional threshold, and most
answer-stage comparisons had no discordant pair.

\subsection{Independent and Dependent Variables}

The independent variables are the three factors above. Retrieval quality was
measured by anchor coverage@5 and MRR@5, scored against gold anchors, and answer
behavior by a ground-truth case crossed with a mechanically derived outcome.
Anchor coverage@5 is deliberately not called Recall@5: it is the fraction of a
question's own gold anchors found in the top \thesisTopK{}, and the standard
name would invite the standard semantics.

\section{Corpus}

The corpus came from \thesisCorpusDataset{} \parencite{kurry2025}, S\&P~500
earnings call transcripts under the MIT License, at revision
\thesisCorpusRevision{}. It covers \thesisNumCompanies{} large-capitalization
United States banks over \thesisNumQuarters{} quarters, 2023 Q1 to 2024 Q4:
\thesisNumTranscripts{} transcripts, every company present throughout.

The sector was chosen for the property the design stresses. Competitors reporting
on the same quarters discuss the same topics in near-identical language, so
telling a passage from a plausible wrong one turns on company and period. Where
documents differ in subject matter, content terms alone tell them apart.

The corpus was frozen before benchmark annotation began. Gold anchors are
verbatim substrings of the transcripts, so an upstream revision leaving the
identifier list intact would invalidate every anchor silently. A digest check
guards against this (appendix).

\section{Preprocessing and Chunking}

Each transcript was divided into two sections, prepared remarks and
question-and-answer, from the per-turn speaker array.
A heuristic located the boundary from the operator's conventional phrasing; no
structural marker was used. The rule was checked across all
\thesisNumTranscripts{} transcripts; that check is documented in the
implementation rather than recorded as an artifact. Each chunk inherits the
section, speaker, and transcript metadata of its record.

Chunks were built as non-overlapping fixed-size token windows, at
\thesisChunkSizeSmall{} and \thesisChunkSizeLarge{} tokens. Token counts were
measured with the embedding model's own tokenizer, that of
\texttt{BAAI/bge-small-en-v1.5} \parencite{baai2023}, so that a chunk's nominal
length matches what the encoder actually consumes. Windows were cut within a
section, so no chunk spans the boundary between prepared remarks and
question-and-answer.

Chunk boundaries are identical across the two indexing representations; only the
string handed to the retrievers differed. Any difference between the
representations is therefore attributable to what was indexed. A chunk identifier refers to the same text only within
one chunk size.

\section{Indexing Representation}

\subsection{The Two Representations}

The third factor controls what string the retrievers index. Under text-only
indexing, the indexed string is the raw chunk text; under metadata-enriched
indexing, it is a single metadata line (the prefix), a blank line, and then that
same chunk text. The line was generated from one template:

\begin{quote}
\ttfamily
\{company\} (\{ticker\}) \{quarter\_prose\} \{year\} \\
Q\{quarter\} \{year\} \{section\}
\end{quote}

\noindent
The template was fixed before any enriched retrieval ran (appendix).

\subsection{Invariants}

Anchor matching always runs against the raw chunk text and never against the
enriched string, so template text could not count as evidence. Both retrievers
receive the identical string, so indexing representation is not confounded with
retrieval strategy. The enriched string serves indexing only and never reaches
the generator.

A \thesisChunkSizeLarge{}-token chunk plus the prefix overruns the encoder's
content budget, so the text was truncated when the index was built, keeping
both retrievers on one string (appendix). Truncation can make a chunk harder to retrieve but
does not change how it is scored, since anchor matching runs against the raw
chunk text.

\section{Retrieval}

Dense retrieval embedded chunks and queries with
\texttt{BAAI/bge-small-en-v1.5} \parencite{baai2023} and stored the vectors in a
FAISS index. Queries carried the instruction prefix the model's authors document
\parencite{baai2023}, \texttt{Represent this sentence for searching relevant
passages:}. Passages were embedded without it.

Sparse retrieval used BM25 \parencite{robertson2009}, through the \texttt{BM25Okapi}
implementation of \texttt{rank\_bm25}, over the same indexed string. The
tokenizer lowercased the text and split it on alphanumeric runs, so every ticker
survived as a token.

The hybrid strategy combined the two ranked lists by reciprocal rank fusion
\parencite{cormack2009}. It scores a chunk by summing, over both rankings,
the reciprocal of \thesisRrfK{} plus its rank. Each retriever ranked the
entire chunk set, so fusion saw both complete orderings before the fused list
was cut to \thesisTopK{}. Fusion operates on ranks, so the retrievers' score
scales are never compared.

Retrieval ran against the whole corpus. Every chunk of all
\thesisNumTranscripts{} transcripts was a candidate for every question.
Retrieval quality is therefore measured on a
deliberately hard task, so the absolute figures support comparison between
conditions, not an absolute estimate.

\section{Benchmark}

\subsection{Composition and Authorship}

The benchmark holds \thesisNumQuestions{} questions carrying
\thesisNumGoldAnchors{} gold anchors. Of these questions,
\thesisNumQuestionsFactual{} are factual, \thesisNumQuestionsThematic{} thematic, \thesisNumQuestionsComparative{}
comparative, and \thesisNumQuestionsUnanswerable{} unanswerable; the unanswerable
questions carry none.

For the unanswerable questions, absence of an answer was proven:
their search terms were re-counted at build time, and one planned question was
dropped when its bank's name appeared 15 times.

Candidates were generated by a large language model
operating over the frozen corpus under progressively formalized rule
constraints, with each iteration reviewed by the researcher and the defects that
review identified driving the next rule. The researcher did not author each
question independently from a first reading of the transcripts. What the
researcher authored is the rule set, the review at each iteration, and the
accept-or-reject decision on every item and every anchor. No model decided what
entered the benchmark.

\subsection{Gold Anchors Rather Than Chunk Identifiers}

A gold anchor is a short verbatim string lifted from the source transcript. A
retrieved chunk counts as a hit when it contains that string as a substring.
Chunk identifiers were rejected because an identifier is not stable across chunk
size: the 47th chunk at one size is different text from the 47th at the other.

In eight items, an anchor's margin left an essential term outside any valid
span, and the question was permitted to name it. This raised question-to-passage
lexical overlap in those items and may have favored the lexical retrievers.

\section{Measures}

\subsection{Retrieval Quality}

Write $A_q$ for the gold anchors of question $q$ and $\mathrm{rank}(a)$ for the
one-indexed position of the first of the top $k = \thesisTopK{}$ retrieved
chunks containing anchor $a$ verbatim, undefined when no chunk does. An anchor's
presence is checked against the raw chunk text, never the indexed string, which
under enrichment carries a metadata line the transcript does not.

\[
  \mathrm{coverage@5}(q)
  = \frac{\bigl|\{ a \in A_q : \mathrm{rank}(a) \text{ defined} \}\bigr|}
         {\lvert A_q \rvert}
  \qquad
  \mathrm{MRR@5}(q) = \frac{1}{\lvert A_q \rvert} \sum_{a \in A_q} c(a)
\]

\noindent
with $c(a) = 1/\mathrm{rank}(a)$ when that rank is defined and zero otherwise.
Partial retrieval therefore earns partial credit but never full credit.

MRR@5 keeps a standard name for a non-standard quantity. Reciprocal rank is
conventionally averaged \emph{across queries}, one per query. Here the inner
mean is taken \emph{within} a question, across its own anchors. The two
definitions coincide only when a question has one anchor.

Unanswerable items carry no anchors and are excluded rather than scored zero, so
every coverage and rank figure has \thesisNumQuestionsScored{} as its
denominator, not \thesisNumQuestions{}.

\subsection{Answer Classification}

Answers were classified on two axes: a ground-truth case, set by retrieval, and
an outcome, set by the generator's response. Outcomes are reported per case. Case~A holds the \thesisNumQuestionsUnanswerable{} unanswerable items.
Case~B holds answerable questions whose coverage is zero, and Case~C answerable
questions whose coverage is above zero, where at least one anchor reached the
context.

\emph{Abstained} is assigned when the canonical refusal sentence appears
anywhere in the response as a complete sentence, matched byte for byte and
case-sensitively. Otherwise, every numeric claim is checked against the
retrieved text. The response is \emph{answered, grounded, on target} when every
claim appears there verbatim and a supporting chunk belongs to the question's
own transcripts. It is \emph{answered, grounded, off target} when no supporting
chunk belongs to them, and \emph{answered, ungrounded} when any claim traces to
nothing retrieved. A Case~C record that is answered, grounded, and on target is said to
have \emph{converted}.

The labels do not show two steps of the grounding check. A bare year or quarter the
question itself supplies is excluded before checking. Where an answer makes no
numeric claim, grounding instead requires a run of six consecutive words shared
with some retrieved chunk, and an answer with no such run counts as ungrounded.

Grounding is a verbatim-substring proxy, likelier to flag a faithful answer than
to pass a fabricated one. The off-target label describes what retrieval
returned, not answer quality.

\subsection{Manual Faithfulness Review}

A packet of \thesisNumReviewItems{} records was drawn from the
\thesisNumGenerations{} generations by a seeded mechanical rule, seed 20260906,
over five disjoint strata. The records were read in a blind, shuffled order. One
stratum is a census of every ungrounded record, so a finding about those records
describes their whole population. The other strata are samples, and a
proportion within them is a rate within the sample. A single reviewer read every
record, so no inter-rater statistic exists.

\section{Generation}

\subsection{Model and Decoding}

Every answer this work reports was produced by \thesisQeightModel{} \parencite{qwen3.8-27b}, a
27-billion-parameter dense model run locally through LM Studio. The generator is
a held-constant factor: one model, one build, and one set of decoding settings
produced all \thesisQeightNumRecords{} records. The build is identified by
weight file and repository \parencite{unsloth-qwen3.8-27b-gguf}, but no digest
of those weights is recorded. A reader can therefore fetch the build but cannot
verify that it is the file that ran, which is weaker provenance than the corpus
carries.

Each call was stateless. Three decoding settings were sent with every request:
temperature 0, seed 20260831, and reasoning effort \texttt{none}. They were
confirmed from a capture of the request body actually sent. The key and value caches were quantized to \texttt{q8\_0}.

A further generator arm exists, a full run of the grid under a different
generator, and it is not a result: \thesisFlashRobustnessModel{}
\parencite{qwen3.8-flash-next} is a robustness check, reported as such in the
chapters that follow.

\subsection{Prompt and Abstention}

The system prompt was fixed before any condition ran and is reproduced in full,
with the format of the user turn (appendix). It
instructs the model to answer only from the retrieved context and, where the
context is insufficient, to respond with one exact canonical sentence and
nothing else. Every record carries a hash of the prompt that produced it, and
all \thesisQeightNumRecords{} records carry the same hash.

One record in \thesisQeightNumRecords{} states the canonical sentence after a
connective, so the case-sensitive test does not fire, and the record is labeled
a Case~B off-target answer. The defect was found after the results existed and
is reported rather than corrected.

\section{Analysis}

\subsection{Descriptive Statistics}

The unit of observation is a question under one condition, which gives 480
observations: \thesisNumQuestionsScored{} questions in each of
\thesisNumConditions{} conditions. Ten variables enter: the design factors and
question category as indicators, the anchor count, chunk-size sensitivity, and
the two outcomes. Coefficients are Pearson throughout: point-biserial for one
dichotomous variable against a continuous one, phi for two dichotomous.

{\centering\noindent Insert Table \ref{tab:descriptives} about here\par}

\subsection{Inferential Tests}

Every test is paired: each question ran through every condition, so each
question's own difficulty cancels out of the comparison.

\subsubsection{Stage 1}

Retrieval quality was compared with the two-sided paired Wilcoxon signed-rank
test \parencite{wilcoxon1945}. A comparison in which every question tied is
untestable and consumes no Holm rank. The effect size is
$r = z / \sqrt{n_{\neq 0}}$ over non-tied pairs, reported only when there are
10 or more; its direction comes from the win and loss counts (appendix).

Holm correction \parencite{holm1979} applies within each family. Family 1
compares each metadata-enriched condition with its matched text-only
counterpart, six pairs in all. Families 2 and 3 hold all pairwise comparisons
within each representation, 15 pairs each. Each family is corrected separately
for each metric, so a claim that names no metric is wrong about at least one of
them. Question categories form separate families.

\subsubsection{Stage 2}

Generation outcomes are binary at the level of a matched pair, so they were
compared with McNemar's test \parencite{mcnemar1947} in its \emph{exact} binomial
form. This is a two-sided binomial test on the larger discordant count against
$n = b + c$ trials at a null probability of .50. With four discordant pairs the
smallest attainable \emph{p} value is .125, so a floor that recurs in the
results reflects the test's granularity.

The unit of analysis is the question. A question contributes up to
\thesisNumConditions{} observations that are not independent, so an exact
binomial over observations would reject the null too readily. Collapsing to one
observation per question removes the dependence, and the question-level test is
the one quoted.

Where two conditions produced identical outcomes on every question, there are no
discordant pairs and no \emph{p} value exists. This held for most comparisons
across the full answer-stage grid, and the testable count is given beside the
total wherever such comparisons are reported.

One convention governs every null result reported in this thesis:

\begin{quote}
A test that does not reject the null does not by itself bound the size of an
effect that may exist here, because no confidence interval accompanies these
tests and, at the informative-pair counts these data commonly produce, the
smallest attainable \emph{p} value can exceed any conventional threshold
whatever the effect. Such a result is reported as a failure to detect a
difference at this sample size, never as evidence that the conditions are
equivalent, and every statement of one carries the sample size and the power
limitation that produced it.
\end{quote}

An informative pair is one that can move the test: for McNemar the
discordant pairs, for the signed-rank test the pairs whose difference is not
zero.
Unlike nearly all pooled retrieval comparisons, many within one question
category have an uncorrected smallest attainable \emph{p} value above .05 and
could not reject whatever the effect; corrected attainability depends on the
family, and the analysis record lists each.

The research questions, the corpus, the benchmark, the retrieval measures, and
two of the three manipulated factors, chunk size and retrieval strategy, were
fixed before any condition was run. The third factor, indexing representation,
was added after the six conditions that vary only the first two had been run,
which is why those six form one half of the grid unchanged. No part of this
work was pre-registered. The pipeline is logged, pinned, and re-runnable
(appendix).

Finally, where an analysis in this work other than the tests of the three
hypotheses was not pre-specified, it is reported as descriptive rather than
confirmatory, labeled as such at the point of use, and no hypothesis is
accepted or rejected on it. Checks of the corpus, the index, and the instrument, including the manual
review's reading of the outcome labels, are not analyses in this sense and fall
outside the rule.
